\chapter{Heegner meets Ramanujan}
\begin{minipage}[t]{\textwidth}
\lettrine{R}{ajeev} Khan sat at his Dehli office desk, a single lamp casting a warm glow over his notes. 
\begin{figure}[H]
\includegraphics[width=1.0\linewidth]{Illustrations/vignetter-reggi.jpg}
\caption{Rajeev pondering Heegner Numbers}
\end{figure}
\noindent
He was busy ceding his eminence to the quadratic \(\mathbb{F}\)ield extenders of Heegner primes, \(
\{2, 3, 5, 7, 11, 19, 43, 67, 163\}
\).
\end{minipage}

Each was a node in a web of transcendence and regularity, their properties spiraling outward into modular forms, elliptic curves, and the near-integrality of the Ramanujan number:
\[
e^{\pi \sqrt{163}}.
\]

He paused, still haunted by the equation. Why this particular arrangement with \(\pi\) and \(e\)? The answers, he knew, lay somewhere just a after mathematics and before metaphysics.

\subsection*{Why \(\pi\) Multiplies \(e\)?}

Rajeev began with the heart of the question. The Ramanujan number, \(e^{\pi \sqrt{163}}\), was built from two transcendental numbers: \(e\) and \(\pi\), paired with the square root of a quadratic number. But why should \(\pi\) serve as the multiplier while \(e\) acted as the base?

\begin{quote}
"It’s the geometry," he thought. \(\pi\) governs circles, periodicity, and symmetry. It’s the constant that ties the circumference of a circle to its diameter, the multiplier that embeds periodicity into the fabric of space. Meanwhile, \(e\) governs exponential growth and decay, the natural rhythm of change and continuity.
\end{quote}

Together, they described a world that was at once cyclical and expansive. No less a dance between repetition and transcendence. In the Ramanujan number, \(\pi\) as the multiplier brought periodicity to the exponential growth of \(e\), anchoring it within the modular structure of the quadratic field. Without \(\pi\), the connection to modular forms, which are the symmetries that arise from the lattice structure of elliptic curves, would unravel. The result would no longer reflect the profound near-integer property that makes \(e^{\pi \sqrt{163}}\) so extraordinary.

\begin{quote}
"The multiplier must be \(\pi\)," he wrote, "because it embeds circularity into the exponential, weaving periodicity into the infinite."
\end{quote}

The Heegner primes themselves were remarkable not just for their connection to unique factorization, but for their regularity. Rajeev flipped through his notes, confirming what he suspected: every single Heegner prime was regular. None of them divided the numerators of the critical Bernoulli numbers, \(B_2, B_4, \dots, B_{p-3}\). Such regularity was not so special among primes.

\begin{quote}
"And they must be primes," he thought. "Only primes carry the purity needed to define the arithmetic of these fields."
\end{quote}

\begin{figure}[H]
\includegraphics[width=1.0\linewidth]{Illustrations/vign-rajeev2.png}
\caption{Rajeev before taking his post at the library}
\end{figure}

The necessity of primality tied back to the Heegner property: these numbers defined imaginary quadratic fields with class number 1, where unique factorization held. If any of the Heegner numbers were composite, their corresponding quadratic fields would lose this property. The dance of modular symmetry would collapse.
Proffering more of the same he muttered:
\begin{quote}
"The regularity of the Heegner primes ensures their stability within the modular forms. They anchor the infinite expansions of \(e^{\pi \sqrt{d}}\) in the finite world of integers."
\end{quote}

\subsection*{A Harmony Between Integer and Infinite}

The Heegner primes, through their innate regularity, encoded the most harmonious of quadratic fields. The Ramanujan number captured that harmony in the interplay of the transcendental and quadratic.

\begin{quote}
"It’s like the universe whispers in integers," Rajeev mused. "But those whispers are wrapped in infinite expansions of:

\(e, \pi, \text{ and } \sqrt{d}.\)"
\end{quote}

The near-integrality of \(e^{\pi \sqrt{163}}\) arises because 163 was the crown jewel of quadratic fields. The exponential and circular properties of \(e\) and \(\pi\) aligned perfectly with the modular symmetries of the field, creating a number that hovered on the edge of discreteness. Rajeev pondered:

\begin{quote}
"Without primality, the modular forms would lose their anchor," he wrote. "The unique factorization of the quadratic fields depends on the purity of primes. Only primes can stabilize the infinite expansions of \(e^{\pi \sqrt{d}}\)."
\end{quote}

\begin{figure}[H]
\includegraphics[width=1.0\linewidth]{Illustrations/vign-rajeevD.png}
\caption{Rajeev still pondering Heegner Numbers}
\end{figure}

Primality, regularity, modularity: these were all aspects of the same phenomenon. Rajeev staring at his notebook, 

\begin{quote}
"The infinite reaches of \(e\), the symmetry of \(\pi\), and the stability of primes—they remind us that the infinite and the finite are not so distinct."
\end{quote}

With that, he closed his notebook and pondered the expansion forms for e and \(pi\) that would suffice to code for in order to deliver essentially an integer Ramanujan number.

% >>>>>> ANNEX E (cut-sheet v2): Euler's totient, the group U_n and the RSA key algorithm — block commented out below, pointer left in text <<<<<<
\textit{(The machinery itself --- totients, the group $U_n$, and the RSA recipe built on them --- is laid out in Annexe E.)}

% \section*{Euler's Totient Function $\phi(n)$}
% 
% Since $n = 12 = 2^2 \times 3$, we count numbers less than 12 that are coprime to it. Thus, given greatest common divisor (gcd):
% \begin{align*}
% \gcd(1,12)&=1, \quad \gcd(5,12)=1, \quad \gcd(7,12)=1, \quad \gcd(11,12)=1.
% \end{align*}
% \noindent
% More generally, for a prime number $p$, we have $\phi(p) = p-1$, as all numbers below it are coprime and for composite $n$ with distinct prime factors:
% \[
% \phi(n) = n\prod_{p \mid n} \left(1 - \frac{1}{p}\right).
% \]
% Applying this to $\phi(12)$:
% \[
% \phi(12) = 12 \left(1 - \frac{1}{2}\right)\left(1 - \frac{1}{3}\right) = 12 \times \frac{1}{2} \times \frac{2}{3} = 4.
% \]
% The Inclusion-Exclusion principle corrects for overcounting in sets. Given two sets $A$ and $B$:
% \(
% |A \cup B| = |A| + |B| - |A \cap B|.
% \) Applying this to $\phi(54)$, where $54 = 2 \times 3^3$:
% \[
% \phi(54) = 54 \left(1 - \frac{1}{2} - \frac{1}{3} + \frac{1}{6}\right) = 54 \times \frac{1}{3} = 18.
% \]
% For a prime power $p^j$ as multiples of $p$ up to $p^j$ are not coprime:
% \(
% \phi(p^j) = p^j - p^{j-1}.
% \)
% The dot plot generated via \href{https://colab.research.google.com/drive/1-TYYGCmru1JCt1gVpfZrCToPLPpUNtTV?usp=sharing}{GraphEulerTotient6n+1.ipynb} displays coprimes as dots and $\phi(n)$ as a line graph.
% \begin{figure}[H]
%     \centering
%     \includegraphics[width=0.8\textwidth]{Illustrations/totient-function.png}
%     \caption{Euler's totient function with coprime dot plot.}
%     \label{totientFunction:fig}
% \end{figure}
% \newpage
% \section*{$U_n$: The Group of Invertible Integers Modulo $n$}
% 
% Euler's totient function, $\phi(n)$ helps define the group of invertible integers modulo $n$, denoted as $U_n$. The set $U_n$ consists of integers less than $n$ that are coprime to $n$, and it forms a multiplicative group under modulo $n$ multiplication:
% \begin{equation}
% U_n = { a \in \mathbb{Z} \mid 1 \leq a < n, \gcd(a, n) = 1 }.
% \end{equation}
% The number of elements in $U_n$ is exactly $\phi(n)$. If $n$ is prime, then every integer from 1 to $n-1$ is coprime to $n$, so $\phi(n) = n - 1$. An element $a \in U_n$ is called a unit involution if:
% \begin{equation}
% a^2 \equiv 1 \pmod{n}.
% \end{equation}
% That is, each element is its own inverse. For example, in $U_8 = {1, 3, 5, 7}$:
% \begin{itemize}
% \item $1^2 \equiv 1 \mod 8 \quad 3^2 \equiv 9 \equiv 1 \mod 8$
% \item $5^2 \equiv 25 \equiv 1 \mod 8$ \quad $7^2 \equiv 49 \equiv 1 \mod 8$
% \end{itemize}
% Thus, all elements of $U_8$ are unit involutions.
% 
% The {\em order} of the group $U_n$ is given by $\phi(n)$, the {\em totient} function of $n$, because the number of elements in $U_n$ is exactly the count of numbers less than $n$ that are coprime to $n$. When \(n\) is a prime number \(p\), the order of the group \(U_p\) is given as \( \phi(p) = p-1 \) since for a prime number \(p\), every integer from 1 to \(p-1\) is relatively prime\footnote{since none of those integers share any common factors with \(p\) other than 1. For a prime number \(p\), we simply count all numbers from 1 up to \(p-1\), as all of these numbers do not share any divisors with \(p\).} to \(p\).  
% 
% \begin{figure}[h!]
% \centering
% \includegraphics[width=0.6\textwidth]{Illustrations/invertibles.png}
% \caption{Group of Invertibles }
% \end{figure}
% 
% \subsection*{Role of $U_n$ in RSA}
% The RSA cryptosystem relies on the structure of $U_n$, where $n = pq$ is a product of two large prime numbers $p$ and $q$. Euler's totient function is crucial in RSA as:
% \begin{equation}
% \phi(n) = (p-1)(q-1).
% \end{equation}
% The public key exponent $e$ is chosen such that:
% \begin{equation}
% e \in U_{\phi(n)},
% \end{equation}
% meaning that $e$ is coprime to $\phi(n)$ and has a modular inverse. The private key exponent $d$ is the modular inverse of $e$ modulo $\phi(n)$:
% \begin{equation}
% ed \equiv 1 \mod \phi(n).
% \end{equation}
% RSA encryption follows three steps:
% \begin{enumerate}
% \item \textbf{Key Generation} - recipient generates public $(n, e)$ and private key $(d, n)$.
% \item \textbf{Encryption} - sender encrypts a message using the public key.
% \item \textbf{Decryption} - recipient decrypts the message using the private key.
% \end{enumerate}
% 
% \subsection*{Key Generation Algorithm}
% \small
% \begin{lstlisting}[language=Python, basicstyle=\ttfamily\small, frame=single]
% def generate_keys():
% p = randprime(103, 104)
% while q == p:
% q = randprime(103, 104)
% n = p * q
% phi_n = (p - 1) * (q - 1)
% e = random.randrange(2, phi_n)
% while gcd(e, phi_n) != 1:
% e = random.randrange(2, phi_n)
% d = mod_inverse(e, phi_n)
% return (n, e), (d, n)
% \end{lstlisting}
% 
% \subsection*{RSA Encryption}
% \small
% \begin{lstlisting}[language=Python, basicstyle=\ttfamily\small, frame=single]
% def encrypt(public_key, plaintext):
%     n, e = public_key
%     ciphertext = [pow(ord(char) - ord('A') + 1, e, n) 
%                   for char in plaintext.upper() if 'A' <= char <= 'Z']
%     return ciphertext
% \end{lstlisting}
% 
% \subsection*{RSA Decryption}
% \small
% \begin{lstlisting}[language=Python, basicstyle=\ttfamily\small, frame=single]
% def decrypt(private_key, ciphertext):
%     d, n = private_key
%     plaintext = [chr((pow(char, d, n) - 1) % 26 + ord('A')) for char in ciphertext]
%     return ''.join(plaintext)
% \end{lstlisting}
% 
% \newpage
% >>>>>> end of annexed block <<<<<<
\section*{A Conversation on Injective Yet Obscure Mappings: \newline Rajeev and Eleanor Discuss Cryptography}

Rajeev leaned back in his chair, fingers tracing the rim of his coffee cup, eyes glinting with amusement. \textit{“You ever wonder,”} he mused, \textit{“how to build a function that looks friendly but hides its teeth?”}
\bigskip
Eleanor raised an eyebrow, setting down her pen. \textit{“A trapdoor function?”}

\bigskip

\textit{“Not just any trapdoor function. One that’s deterministic but still impossible to reverse.”} He gestured idly toward his notebook, filled with scribbles of number-theoretic expressions. \textit{“No pseudo-random hacks. Just structure, pure and unbreakable.”}

\bigskip
\noindent
Eleanor took a sip of tea, nodding. \textit{“Well, let’s start simple. You want an injective function?”}
\bigskip

\noindent
Rajeev tapped his notebook. \textit{“Consider the Euler totient function: \( T(n) = \sum_{k=1}^{n} \phi(k) \). Strictly increasing, always unique. Easy to invert with binary search.”}
\bigskip

\noindent
Eleanor smiled wryly. \textit{“That’s your problem. If I can invert it quickly, it’s no good for cryptography. We need more chaos.”}

\bigskip

\noindent
Rajeev nodded. \textit{“That’s why I refined it. Instead of summing over \( n \), sum over primes}:

\[ H(n) = \sum_{p_k \leq P_n} \phi(p_k) \]

Eleanor tilted her head, considering. \textit{“So now inversion depends on the prime counting function \( \pi(x) \). That’s better.”} She reached for his notebook, flipping through the pages. \textit{“But it’s still possible, given enough computational power. We need something more obscure.”}

\bigskip

Rajeev grinned, as though she had walked straight into his trap. \textit{“Exactly. That’s where I add modular transformations. Define \( M(n) = H(n) \mod m_n \), where \( m_n \) is a prime-based modulus.”}

Eleanor’s eyes flickered with interest. \textit{“You’re disrupting monotonicity. Nice. Now inversion requires knowing both \( H(n) \) and the sequence \( m_n \).”}
\bigskip

Rajeev on it now. \textit{“Better yet, \( m_n \) isn’t random—it’s structured. Something like \( p_{f(n)} \), where \( f(n) \) is a non-trivial mapping of \( n \).”}

Eleanor tapped her chin affectedly. \textit{“So to invert \( M(n) \), you’d need to untangle both the prime distribution and the modular sequence. That’s a serious roadblock.”}

\bigskip

Rajeev chuckled. \textit{“And yet, it’s not enough. True trapdoor functions need inversion to be infeasible even if you know the structure.”}

Eleanor, pretending to idly watch the afternoon light filter through the window. \textit{“You sound like someone who’s had a revelation.”}
\bigskip

Rajeev shrugged. \textit{“Let’s just say I have a renewed appreciation for randomness. But if we can crack this problem, Eleanor, we might just change the landscape of cryptography.”}

Eleanor smirked. \textit{“Then let’s see how far we can push it.”}

\bigskip

\section*{Number Amicability}

The concept of amicable numbers has evolved through historical definitions.

A trapdoor function is a function that is easy to compute in one direction but difficult to invert without some secret knowledge. In the annex we explore the difficulty of finding such functions by considering the structure of amicable numbers, the Euler totient function, and its cumulative sum.

% >>>>>> ANNEX E (cut-sheet v2): Euclid's and Euler's perfect-number extensions and why inversion is hard — block commented out below, pointer left in text <<<<<<
\textit{(Euclid's and Euler's perfect-number extensions and why inversion is hard --- the full working is set out in Annexe E.)}

% \subsection*{Euclid's Definition}
% Euclid defined amicable numbers in terms of the sum of their 	extbf{proper divisors} (excluding the number itself). Two numbers $m$ and $n$ are amicable if:
% \begin{equation*}
% s(m) = n, \quad s(n) = m
% \end{equation*}
% where $s(x)$ is the sum of proper divisors of $x$.
% 
% \subsection*{Euler's Extension}
% Euler extended Euclid’s definition by introducing the function $\sigma(n)$, which sums 	extbf{all} divisors of $n$, including $n$ itself:
% \begin{equation*}
% s(n) = \sigma(n) - n.
% \end{equation*}
% Rewriting Euclid's definition in terms of $\sigma(n)$, we obtain:
% \begin{align*}
% \sigma(m) - m &= n, \quad \sigma(n) - n = m.
% \end{align*}
% Adding these equations and rearranging, we get:
% \begin{equation*}
% \sigma(m) + \sigma(n) = 2(m + n).
% \end{equation*}
% Since both $m$ and $n$ contribute equally to the sum, the necessary and sufficient condition for $m$ and $n$ to be amicable numbers is:
% \begin{equation}
% \sigma(m) = m + n = \sigma(n).
% \end{equation}
% 
% This extension highlights the broader implications of including all divisors, such as its role in divisor product identities and multiplicative number theory.
% 
% The smallest amicable pairs are:
% \begin{itemize}
% \item \textbf{(220, 284)}
% \begin{align*}
% \sigma(220) &= 1+2+4+5+10+11+20+22+44\\ &+55+110+220 = 284 \\
% \sigma(284) &= 1+2+4+71+142+284 = 220.
% \end{align*}
% \end{itemize}
% 
% Since $\sigma(n)$ encodes divisor structure, one might consider whether a related function such as the \textbf{Euler totient cumulative sum} could exhibit similar \href{https://colab.research.google.com/drive/1GAwbrSXMH0O-KXmb3OMSJwsRj7uwMqg5?usp=sharing}{properties}.
% The fundamental requirement for Amicable numbers $(m, n)$ to satisfy can be rewritten in terms of the Euler totient function $\phi(n)$, which counts the number of integers coprime to $n$:
% 
% \begin{equation*}
%     \sum_{d \mid m} d - m = n, \quad \sum_{d \mid n} d - n = m.
% \end{equation*}
% 
% Define the cumulative totient sum:
% \begin{equation}
%     T(n) = \sum_{k=1}^{n} \phi(k).
% \end{equation}
% 
% This function forms a strictly increasing sequence, creating a lattice structure in the integer domain. Unlike amicable numbers, where a relationship between two numbers is established through divisor sums, $T(n)$ provides a direct summation-based growth pattern.
% 
% \bigskip
% 
% Despite its nontrivial structure, the function $T(n)$ is still \textbf{efficiently invertible}. Given a random number $T_r$ in the sequence, we can determine the corresponding $n$ using binary search in $O(\log n)$ time. Since $T(n)$ is strictly increasing, the mapping $n \to T(n)$ is one-to-one, ensuring that an inverse exists and can be computed efficiently.
% 
% \bigskip
% 
% For a function to be a valid cryptographic trapdoor function, its inverse must be **computationally infeasible** to compute without special knowledge. The totient cumulative sum provides an example of a function that might appear difficult to invert, but in reality, is efficiently computable due to its monotonicity and the ability to apply binary search.
% \bigskip
% 
% This raises the challenge in designing true trapdoor functions: while many number-theoretic functions grow nonlinearly, their invertibility is often easier than expected due to structured behavior or monotonicity.
% \bigskip
% 
% While the sum-of-divisors function used in defining amicable numbers suggests a complex structure, and the cumulative totient sum creates an interesting integer lattice, neither function possesses the necessary computational asymmetry required for a genuine trapdoor function. The difficulty in finding such functions lies in ensuring that inversion remains intractable **even with precomputed knowledge**, an ongoing challenge in cryptographic function design.
% 
%%%%%%%%%%%%
% \section*{Constructing Injective Yet Obscure Mappings: A Pedagogical Path to Trapdoor Functions}
% 
% A cryptographic trapdoor function must be:
% \begin{itemize}
%     \item \textbf{Injective} (one-to-one): Each input must map uniquely to an output.
%     \item \textbf{Computationally asymmetric}: The forward computation should be efficient, while inversion should be infeasible.
% \end{itemize}
% 
% Constructing such a function is difficult without introducing pseudo-random elements. Rajeev was exploring the stepwise evolution of a mapping, moving from simple cumulative sums to prime-based transformations and finally to modular obfuscation, demonstrating the increasing complexity in preventing easy inversion.
% 
% A natural starting point is the cumulative sum of a number-theoretic function. Consider the Euler totient function:
% 
% \begin{equation}
%     T(n) = \sum_{k=1}^{n} \phi(k).
% \end{equation}
% 
% Since \( \phi(k) \) is always positive, \( T(n) \) is **strictly increasing**, making it:
% \begin{itemize}
%     \item \textbf{Injective}: No two distinct \( n \) values share the same \( T(n) \).
%     \item \textbf{Easily Invertible}: Given \( T_r \), we can recover \( n \) via binary search in \( O(\log n) \).
% \end{itemize}
% 
% Thus, \( T(n) \) fails as a trapdoor function because **its bijectivity is immediately obvious**.
% 
% \section{Step 2: Prime-Indexed Totient Encoding (PITE)}
% 
% To disrupt simple inversion, we replace the cumulative sum over natural numbers with one based on \textbf{prime indices}:
% 
% \begin{equation}
%     H(n) = \sum_{p_k \leq P_n} \phi(p_k),
% \end{equation}
% 
% where \( p_k \) is the \( k \)-th prime and \( P_n \) is the \( n \)-th prime.
% 
% \subsection{Why This is Harder to Invert}
% Unlike \( T(n) \), which follows a smooth curve, \( H(n) \) is shaped by:
% \begin{itemize}
%     \item \textbf{Irregular Prime Gaps}: Primes do not follow an arithmetic progression.
%     \item \textbf{Prime-Counting Difficulty}: Estimating \( n \) from \( H(n) \) links inversion to the **prime counting function** \( \pi(x) \), a non-trivial number-theoretic function.
% \end{itemize}
% 
% While \( H(n) \) is still injective, its bijectivity is not immediately apparent, making inversion substantially harder. However, **given enough precomputed data, binary search still works**, so we need an additional layer of obfuscation.
% 
% \section{Step 3: Modular Transformations to Disrupt Bijectivity}
% 
% To further obscure inversion, we introduce a modular transformation:
% 
% \begin{equation}
%     M(n) = H(n) \mod m_n,
% \end{equation}
% 
% where \( m_n \) is a sequence of moduli derived from **structured** rather than random elements, such as:
% 
% \begin{equation}
%     m_n = p_{f(n)},
% \end{equation}
% 
% where \( f(n) \) is a deterministic function mapping \( n \) to an index in the prime sequence.
% 
% \subsection{Why This Further Obscures Inversion}
% \begin{itemize}
%     \item \textbf{Breaks Monotonicity}: Since modulo operations fold values into a bounded range, binary search no longer applies.
%     \item \textbf{Hides Input Structure}: The prime gaps make direct reconstruction of \( n \) from \( M(n) \) infeasible without knowing the transformation sequence.
%     \item \textbf{Remains Injective in a Controlled Domain}: Unlike true randomization, every transformation follows a **deterministic** but **non-trivial** pattern.
% \end{itemize}
% 
% Now, inversion requires **both the original mapping and knowledge of the prime-moduli structure**, making \( M(n) \) a much stronger candidate for a trapdoor function.
% 
% \section{Conclusion: The Challenge of Constructing Trapdoor Functions Without Randomization}
% 
% This progression highlights the difficulty of designing trapdoor functions without relying on pseudo-random elements:
% \begin{enumerate}
%     \item \textbf{Cumulative sums} are **too predictable**, allowing trivial inversion.
%     \item \textbf{Prime-based encodings} disrupt uniform growth, linking inversion to prime distributions.
%     \item \textbf{Modular obfuscation} breaks simple search strategies while preserving injectivity.
% \end{enumerate}
% 
% However, **true trapdoor functions require inversion to remain infeasible even with all known mathematical tools**. Without randomness, ensuring this remains a major open problem.
% 
% \section{The Trade-Off Between Pseudo-Random Elements and Public-Key Protocol Complexity}
% 
% In cryptographic protocols, there is a fundamental trade-off between using pseudo-random elements and maintaining protocol simplicity. While pseudo-randomness strengthens security by introducing unpredictability, it also adds computational overhead and implementation complexity. Avoiding randomness might seem cleaner, but it often exposes the system to structural vulnerabilities.
% 
% To illustrate this trade-off, we examine the Sony PlayStation 3 cryptographic failure, where a deterministic choice in a signature scheme led to the exposure of private keys.
% 
% \section{Case Study: The Sony PlayStation ECDSA Failure}
% 
% Sony's PlayStation 3 used the **Elliptic Curve Digital Signature Algorithm (ECDSA)** for software authentication and encryption. A valid ECDSA signature for a message \( m \) consists of two values \( (r, s) \) computed as follows:
% 
% \begin{enumerate}
%     \item Choose a random integer \( k \) (per signature).
%     \item Compute \( r = (g^k \mod p) \).
%     \item Compute \( s = k^{-1} (H(m) + xr) \mod q \), where \( x \) is the private key.
% \end{enumerate}
% 
% Sony, however, made a crucial mistake: instead of selecting a fresh random \( k \) for each signature, they used a **fixed deterministic value**.
% 
% >>>>>> end of annexed block <<<<<<
\subsection{How This Led to a Security Breach}
Since \( k \) was the same across multiple signatures, an attacker could solve for the private key \( x \) using two signatures \( (r_1, s_1) \) and \( (r_2, s_2) \) on different messages:

\begin{equation}
    s_1 k - s_2 k = H(m_1) - H(m_2) \mod q.
\end{equation}

Rearranging,

\begin{equation}
    k = \frac{H(m_1) - H(m_2)}{s_1 - s_2} \mod q.
\end{equation}

Once \( k \) was known, the private key \( x \) could be computed using:

\begin{equation}
    x = \frac{s_1 k - H(m_1)}{r_1} \mod q.
\end{equation}

This catastrophic failure demonstrated why **deterministic cryptography without randomness is often insecure**.

% \section{Trade-Off: Randomness vs Complexity}

The Sony example illustrates a broader trade-off in cryptographic design:

\subsection{Advantages of Avoiding Randomness}
\begin{itemize}
    \item \textbf{Simplicity:} Deterministic cryptographic operations are easier to implement and verify.
    \item \textbf{Efficiency:} No need for secure random number generation, reducing computation.
    \item \textbf{Predictable Execution:} Avoids timing attacks based on entropy variation.
\end{itemize}

\subsection{Risks of Avoiding Randomness}
\begin{itemize}
    \item \textbf{Precomputability:} Without randomness, attackers can generate lookup tables to break encryption.
    \item \textbf{Replay and Pattern Attacks:} Reusing values in deterministic systems exposes weaknesses, as in the Sony case.
    \item \textbf{Loss of Indistinguishability:} Randomized encryption ensures that the same plaintext never encrypts to the same ciphertext.
\end{itemize}

% \section{Pseudo-Randomness in Public-Key Protocols}

Most modern cryptographic protocols use randomness to prevent structural vulnerabilities:

\begin{itemize}
    \item \textbf{Key Generation:} RSA and ECC require random prime selection to prevent factorization-based attacks.
    \item \textbf{Signature Schemes:} ECDSA and DSA use ephemeral randomness to ensure that private keys remain secure.
    \item \textbf{Encryption:} RSA-OAEP uses randomized padding to prevent chosen-plaintext attacks.
    \item \textbf{Key Exchange:} Diffie-Hellman relies on random exponents to maintain secrecy.
\end{itemize}

While pseudo-randomness increases computational complexity, it significantly improves security by preventing predictable patterns.

\subsection{Can Trapdoor Functions Work Without Randomness?}

If one wishes to avoid randomness, cryptographic security must be derived from inherently hard mathematical problems such as:

\begin{itemize}
    \item \textbf{Integer Factorization (RSA)}
    \item \textbf{Discrete Logarithm Problem (ECDH, DSA)}
    \item \textbf{Lattice-Based Hardness (LWE, NTRU)}
\end{itemize}

However, even these protocols benefit from randomness:
\begin{enumerate}
    \item RSA encryption without randomized padding is vulnerable to chosen-plaintext attacks.
    \item Lattice-based encryption schemes require error terms to prevent efficient inversion.
    \item Quantum-resistant cryptography often introduces random noise to increase computational hardness.
\end{enumerate}

While avoiding randomness simplifies protocol design, it often introduces structural vulnerabilities. As demonstrated by the Sony ECDSA failure, deterministic implementations expose cryptographic keys to attacks.

Thus, randomness is a necessary trade-off:
\begin{itemize}
    \item \textbf{More randomness → Stronger security, but higher computational cost.}
    \item \textbf{Less randomness → Simpler, but vulnerable to structured attacks.}
\end{itemize}

A truly randomness-free trapdoor function would require an alternative complexity source such as structured obfuscation or modular transformations, but these methods introduce their own challenges.

Ultimately, the safest cryptographic systems balance randomness with structured security, ensuring that the function remains computationally asymmetric while preventing predictable vulnerabilities.

\newpage

\section*{Quest for Nine Platonic Solids in Hyperbolic Space }

Celiac sat across from Rajeev Khan in the cozy café, her notebook open to a sketch of a pangolin’s playing with some Platonic solids as they are prone to do. The glow of her tablet illuminated the graph she had been analyzing. Rajeev, sipping his coffee, leaned back, intrigued by her focused energy.

\begin{figure}[H]
\includegraphics[width=1.0\linewidth]{Illustrations/pangolinPlat.png}
\caption{Multipole Moments of Pangolin}
\end{figure}

“You know,” Celiac began, gesturing at the screen, “it was Sylvia who got me thinking about tessellating in hyperbolic space. I was so fixated on positive curvature—like spheres, where \( K > 0 \). But what about the other extreme? Hyperbolic spaces, where \( K < 0 \)?”

Rajeev raised an eyebrow. “You mean spaces with negative curvature? Like a saddle?”

“Precisely,” Celiac said, pointing to her notes. “In flat space (\( K = 0 \)), we’re stuck with the triangle, square, and pentagon as the only allowed polygons for Platonic solids. But in hyperbolic geometry, the situation changes dramatically. The condition we used was:
\[
\frac{1}{p} + \frac{1}{k} < \frac{1}{2} + \frac{1}{\sqrt{-K}},
\]
where \( p \) is the number of sides of the polygons, \( k \) is the number of polygons meeting at a vertex, and \( K < 0 \) is the curvature.”

Rajeev looked at the equation and mused. “And you’re saying this inequality directly determines which \( p \) and \( k \) configurations are allowed in hyperbolic space?”

“Exactly!” Celiac continued. “And there’s more. Each allowed configuration must also satisfy the dihedral angle constraint:
\[
k \cdot \arccos\left(1 - \frac{2}{p}\right) < 2\pi.
\]
By carefully analyzing these conditions across a range of negative \( K \), we discovered something incredible.”

Celiac leaned forward, her voice alive with excitement. “For \( K = 0 \), we already know there are only five Platonic solids. However, by exploring \( K < 0 \), we found that the number of Platonic solids \( S(K) \) increases for certain curvatures, reaching exactly nine at a specific value of \( K \). 

\begin{figure}[h]
    \centering
    \includegraphics[width=0.9\textwidth]{Illustrations/platonicSolids9.png}
    \caption{Number of Platonic Solids (\(S\)) vs. Curvature (\(K\)). The green line marks the exact curvature \(K = -0.4762\) where nine Platonic solids are possible.}
    \label{fig:platonic-solids}
\end{figure}

The code
\href{https://colab.research.google.com/drive/1XIWfotANIzRx8tLbm9sMG_4bei6jpl5l?usp=sharing}{dihedralplatonicHyperbole9.ipynb} 
delivers the figure which shows how the number of Platonic solids \( S(K) \) varies with the curvature \( K \). The red dashed line represents the target \( S(K) = 9 \) and we identified this optimal curvature to be:
\[
K \approx -0.4762 \quad \text{(with fine resolution)}.
\]
Here’s a summary of the result:
\begin{itemize}
    \item In Euclidean flat space (\( K = 0 \)), \( S(K) = 5 \).
    \item For hyperbolic saddle space with \( K \approx -0.4762 \), \( S(K) = 9 \).
\end{itemize}

“Think about it,” Celiac said, her voice alive with curiosity. “If the ancient Greeks had discovered hyperbolic geometry, they might have speculated about saddle spaces filled with nine Platonic solids. Instead, they were limited by their focus on positive curvature and the sphere.”

Rajeev smiled. “You’re saying Sylvia managed to distract you from pangolins to hyperbolic tessellations and gave you a reason to explore a universe with exactly nine Platonic solids?”

Celiac nodded. “And there’s a deeper connection,” she said, pointing to her notes again. “Jovannah once mentioned how she was working on a multipole expansion version of \(\sqrt{d}\) for reciprocal figure sequences other than the squares of \(\zeta(2)\). She did not manage to hit anywhere near the accuracy of the Ramanunjan number but that got me thinking that there are nine Heegner numbers, these can act as coefficients for surface-to-volume ratios of these nine hyperbolic Platonic solids. The alignment is contrived but almost worth exploring.”

Rajeev leaned closer, his brow furrowed. “So you’re saying the nine Heegner numbers somehow correspond to the nine Platonic solids in hyperbolic space?”

“Exactly,” Celiac replied. “It’s as though these mathematical sequences are encoding geometric properties that manifest in hyperbolic spaces. It’s like the universe left us a blueprint.”

Rajeev raised his coffee cup. “To pangolins, hyperbolic geometry, and Heegner numbers. And to the mysteries of exactly nine Platonic solids.”

Celiac laughed, clinking her cup against his. “To curiosity—and the strange beauty of negative curvature.”

% Hyperbolic geometry continues to reveal new insights into the structure of space and the constraints on Platonic solids. By finding a curvature \( K \approx -0.478 \), we uncovered a unique space that hosts exactly nine Platonic solids—a remarkable expansion beyond the familiar five of Euclidean space.







